% Reviewed translation: PDF pages 160--165 / printed pages 146--151.
% The original scan is authoritative; environment headings remain English.
\MathBookSourcePage{160}
由 Lemma~\ref{lem:primitive-tetrahedron-first-unisolvence},
仍可在 $\mathcal{P}_1(\kappa)$ 上定义插值算子 $r_\kappa$:
\[
r_\kappa\mathbf v(a_i)=\mathbf v(a_i),
\qquad
\int_{F_i}(r_\kappa\mathbf v-\mathbf v)\cdot\mathbf n_i\,\mathrm ds=0,
\qquad1\leq i\leq4.
\]
但与二维情形一样, 该算子不满足 Hypothesis H3, 因为它并非定义在
$H^1(\Omega)^3$ 上. 因此, 我们建议用 Bernardi [9] 在
$\mathbb R^3$ 中发展的局部正则化算子的值替代 $\mathbf v$ 的值;
该算子推广了 Section~\ref{sec:appendix-interpolation-discontinuous-functions}
中的二维算子 $R_h$. 这里没有篇幅详细描述这个同样记为 $R_h$ 的算子.
我们只需要知道
$R_h\in\mathcal L(H_0^1(\Omega)^3;\Phi_h^3)$, 且当三角剖分
$\mathcal{T}_h$ 正则时, $R_h$ 满足
\eqref{eq:primitive-bernardi-raugel-local-regularization-error}.
然后定义算子 $\pi_h\in\mathcal L(H_0^1(\Omega)^3;\bar X_h)$:
\begin{equation}\label{eq:primitive-tetrahedron-first-fortin-operator}
\left\{
\begin{aligned}
\pi_h\mathbf v(a)&=R_h\mathbf v(a)
&&\forall\text{ $\mathcal{T}_h$ 的节点 }a,\\
\int_F(\pi_h\mathbf v-\mathbf v)\cdot\mathbf n\,\mathrm ds&=0
&&\forall\text{ $\mathcal{T}_h$ 的面 }F.
\end{aligned}
\right.
\end{equation}
对 Lemma~\ref{lem:primitive-bernardi-raugel-fortin-estimates} 的证明稍作修改,
即可表明 $\pi_h$ 满足 $l=1$ 的 Hypothesis H1, 也满足 Hypothesis H3.

\setcounter{statement}{7}
\begin{Lemma}\label{lem:primitive-tetrahedron-first-fortin-estimates}
由 \eqref{eq:primitive-tetrahedron-first-fortin-operator} 定义的算子
$\pi_h$ 满足
\[
\int_\Omega\operatorname{div}(\mathbf v-\pi_h\mathbf v)q\,\mathrm dx=0
\qquad\forall q\in Q_h.
\]
此外, 若三角剖分 $\mathcal{T}_h$ 正则, 则 $\pi_h$ 满足误差界
\begin{equation}\label{eq:primitive-tetrahedron-first-interpolation-error}
|\mathbf v-\pi_h\mathbf v|_{m,\Omega}
\leq Ch^{k-m}|\mathbf v|_{k,\Omega}
\qquad\forall\mathbf v\in H^k(\Omega)^3,
\end{equation}
其中 $m=0$ 或 $1$, $k=1$ 或 $2$.
\end{Lemma}

最后, 应用 Theorems~\ref{thm:primitive-homogeneous-stokes-convergence} 和
\ref{thm:primitive-stokes-l2-error}, 得到该一阶格式的预期估计.

\setcounter{statement}{3}
\begin{Theorem}\label{thm:primitive-tetrahedron-first-convergence}
设 $\Omega$ 是 $\mathbb R^3$ 中的有界多面体, Stokes 问题的解
$(\mathbf u,p)$ 满足
\[
\mathbf u\in[H^2(\Omega)\cap H_0^1(\Omega)]^3,
\qquad p\in H^1(\Omega)\cap L_0^2(\Omega).
\]
若三角剖分 $\mathcal{T}_h$ 正则, 则
\eqref{eq:primitive-stokes-discrete-mixed} 的解 $(\mathbf u_h,p_h)$,
其中空间 $\bar X_h$ 和 $\bar M_h$ 分别由
\eqref{eq:primitive-tetrahedron-first-velocity-spaces} 和
\eqref{eq:primitive-tetrahedron-first-pressure-spaces} 定义, 满足估计
\begin{equation}\label{eq:primitive-tetrahedron-first-h1-error}
|\mathbf u-\mathbf u_h|_{1,\Omega}
+\|p-p_h\|_{0,\Omega}
\leq C_1h(|\mathbf u|_{2,\Omega}+|p|_{1,\Omega}).
\end{equation}
此外, 若 Stokes 问题
\eqref{eq:primitive-stokes-dual-regularity-problem} 是正则的, 则有
$L^2$ 估计
\begin{equation}\label{eq:primitive-tetrahedron-first-l2-error}
\|\mathbf u-\mathbf u_h\|_{0,\Omega}
\leq C_2h^2(|\mathbf u|_{2,\Omega}+|p|_{1,\Omega}).
\end{equation}
\end{Theorem}

\MathBookSourcePage{161}
现在转向二阶格式. 如所预期, 我们希望构造一个速度向量 $\mathbf w$,
它在 $\kappa$ 的边界上具有二次切向分量, 并与 $\kappa$ 中的仿射压力
相容. 参照 $\mathbb R^2$ 中的相应格式, 人们容易想到在空间
$\{P_2\oplus(\lambda_1\lambda_2\lambda_3\lambda_4P_0)\}^3$
中选取 $\mathbf w$. 遗憾的是, 这个空间的维数太小, 无法满足仿射压力
所要求的 inf-sup 条件. 按照 Fortin [30] 和 Bernardi \& Raugel [10],
最好的办法是在该空间中加入
\eqref{eq:primitive-tetrahedron-first-local-space} 的三次多项式
$\mathbf p_i$; 因此在 $P_4^3$ 的如下子空间中选取 $\mathbf w$:
\begin{equation}\label{eq:primitive-tetrahedron-second-local-space}
\mathcal{P}_2(\kappa)
=\{P_2\oplus(\lambda_1\lambda_2\lambda_3\lambda_4P_0)\}^3
\oplus\operatorname{span}\{\mathbf p_1,\mathbf p_2,\mathbf p_3,\mathbf p_4\}.
\end{equation}
立即可见 $\mathcal{P}_1(\kappa)\subset\mathcal{P}_2(\kappa)$;
因此可以尽可能应用
Subsection~\ref{subsec:primitive-discrete-inf-sup-methods} 的材料来建立
inf-sup 条件. 这意味着 inf-sup 条件不对速度 $\mathbf w$ 的自由度施加
约束, 因而可以选择最方便的自由度. 与 $P_2$ 自然相联系的自由度是
\[
\mathbf w(a_i),\quad1\leq i\leq4,
\qquad
\mathbf w(a_{ij}),\quad a_{ij}\text{ 是 }e_{ij}\text{ 的中点},
\quad1\leq i<j\leq4;
\]
与 $\{\mathbf p_i;1\leq i\leq4\}$ 相联系的是
\[
\int_{F_i}\mathbf w\cdot\mathbf n_i\,\mathrm ds,
\qquad1\leq i\leq4;
\]
而与 ``bubble function'' $\lambda_1\lambda_2\lambda_3\lambda_4$
相联系的最简单自由度是
\[
\mathbf w(a_\kappa),
\qquad a_\kappa\text{ 是 }\kappa\text{ 的中心}.
\]
下面的 Lemma 表明这组矩对 $\mathcal{P}_2(\kappa)$ 唯一可解.

\setcounter{statement}{8}
\begin{Lemma}\label{lem:primitive-tetrahedron-second-unisolvence}
$\mathcal{P}_2(\kappa)$ 中的多项式 $\mathbf p$ 由下列 $37$ 个量唯一确定:
\begin{equation}\label{eq:primitive-tetrahedron-second-dofs}
\left\{
\begin{aligned}
&\mathbf p(a_i),&&1\leq i\leq4,\\
&\mathbf p(a_{ij}),&&1\leq i<j\leq4,\\
&\mathbf p(a_\kappa),& &\\
&\int_{F_i}\mathbf p\cdot\mathbf n_i\,\mathrm ds,
&&1\leq i\leq4.
\end{aligned}
\right.
\end{equation}
此外, $\mathbf p$ 在任意一个面 $F_i$ 上的限制, 只依赖于它在该面上的
自由度.
\end{Lemma}

\begin{Proof}
$\mathcal{P}_2(\kappa)$ 中的多项式有 $37$ 个系数,
这恰好是 \eqref{eq:primitive-tetrahedron-second-dofs} 定义的自由度数.
因此只需证明零矩只能生成零多项式. 在任意一条边 $e_{ij}$ 上,
$\mathbf p$ 的一个分量 $p$ 约化为单变量二次函数. 所以
$p(a_i)=p(a_j)=p(a_{ij})=0$ 蕴含 $p|_{e_{ij}}=0$.
于是, 若 $\mathbf p$ 在任意一个面 $F_i$ 边界上的自由度均消失,
则必有 $\mathbf p|_{F_i}=c\mathbf p_i$; 若再有
$\int_{F_i}\mathbf p\cdot\mathbf n_i\,\mathrm ds=0$, 则 $c=0$.

最后, 若 $\mathbf p|_{\partial\kappa}=0$, 则 $\mathbf p$ 的每个分量
$p$ 都是一个 ``bubble function'', 因而 $p(a_\kappa)=0$ 蕴含
$p=0$ 于 $\kappa$ 中.
\end{Proof}

\MathBookSourcePage{162}
\setcounter{statement}{6}
\begin{Remark}\label{rem:primitive-tetrahedron-second-affine-dofs}
除 $\int_{F_i}\mathbf p\cdot\mathbf n_i\,\mathrm ds$ 外,
\eqref{eq:primitive-tetrahedron-second-dofs} 列出的所有自由度都在仿射变换下
保持. 注意, 由于三次多项式 $\mathbf p_i$ 中出现法向量,
空间 $\mathcal{P}_2(\kappa)$ 本身并不在仿射变换下保持不变.
\end{Remark}

根据以上讨论, 为速度和压力选取如下有限元空间:
\begin{equation}\label{eq:primitive-tetrahedron-second-velocity-spaces}
\left\{
\begin{aligned}
W_h&=\{\,\mathbf w\in\mathcal C^0(\bar\Omega)^3;
\ \mathbf w|_\kappa\in\mathcal P_2(\kappa)
\quad\forall\kappa\in\mathcal T_h\,\},\\
X_h&=W_h\cap H_0^1(\Omega)^3,
\end{aligned}
\right.
\end{equation}
\begin{equation}\label{eq:primitive-tetrahedron-second-pressure-spaces}
\left\{
\begin{aligned}
Q_h&=\{\,q\in L^2(\Omega);
\ q|_\kappa\in P_1\quad\forall\kappa\in\mathcal T_h\,\},\\
M_h&=Q_h\cap L_0^2(\Omega).
\end{aligned}
\right.
\end{equation}
Lemma~\ref{lem:primitive-tetrahedron-second-unisolvence} 提供适当的插值算子
$r_h$:
\[
r_h\mathbf v|_\kappa=r_\kappa\mathbf v
\qquad\forall\kappa\in\mathcal T_h,
\quad\forall\mathbf v\in\mathcal C^0(\bar\Omega)^3,
\]
其中 $r_\kappa\mathbf v$ 是 $\mathcal P_2(\kappa)$ 中与 $\mathbf v$
在 $\kappa$ 上具有相同自由度
\eqref{eq:primitive-tetrahedron-second-dofs} 的唯一多项式. 显然
\[
r_h\in\mathcal L(\mathcal C^0(\bar\Omega)^3;W_h)
\cap\mathcal L([\mathcal C^0(\bar\Omega)\cap H_0^1(\Omega)]^3;X_h),
\]
而下面的 Lemma 验证 Hypothesis H1.

\setcounter{statement}{9}
\begin{Lemma}\label{lem:primitive-tetrahedron-second-interpolation}
若 $\mathcal T_h$ 是 $\bar\Omega$ 的正则三角剖分, 则 $r_h$ 满足逼近性质
\begin{equation}\label{eq:primitive-tetrahedron-second-interpolation}
|\mathbf v-r_h\mathbf v|_{m,\Omega}
\leq Ch^{k-m}|\mathbf v|_{k,\Omega}
\qquad\forall\mathbf v\in H^k(\Omega)^3,
\end{equation}
其中 $m=0,1$, $k=2$ 或 $3$, 正常数 $C$ 与 $h$ 和 $\mathbf v$ 无关.
\end{Lemma}

\begin{Proof}
若 $r_\kappa$ 在仿射变换下保持, 证明将十分简单;
不过, 实际证明仍很短, 因为只有面上的矩在仿射变换下不保持.
用 $I\mathbf v$ 表示由下式定义的
$P_2\oplus(\lambda_1\lambda_2\lambda_3\lambda_4P_0)$ 中的多项式:
\[
\begin{aligned}
I\mathbf v(a_i)&=\mathbf v(a_i),&&1\leq i\leq4,\\
I\mathbf v(a_{ij})&=\mathbf v(a_{ij}),&&1\leq i<j\leq4,\\
I\mathbf v(a_\kappa)&=\mathbf v(a_\kappa).
\end{aligned}
\]
一方面, 这组条件唯一确定 $I\mathbf v$; 另一方面, 算子 $I$ 在仿射变换下
不变. 此外, 由于 $I$ 保持 $P_2$ 中的多项式,
Corollary~\ref{cor:appendix-abstract-interpolation-physical} 的标准论证表明
\begin{equation}\label{eq:primitive-tetrahedron-second-affine-interpolation}
|\mathbf v-I\mathbf v|_{m,\kappa}
\leq C_1\|B_\kappa\|^k\|B_\kappa^{-1}\|^m
|\mathbf v|_{k,\kappa},
\qquad m=0,1,\quad k=2,3.
\end{equation}
与 Remark~\ref{rem:primitive-bernardi-raugel-decomposition} 中一样,
可将 $r_\kappa\mathbf v$ 表示为
\[
r_\kappa\mathbf v
=I\mathbf v+\sum_{i=1}^4\alpha_i\mathbf p_i
+\boldsymbol\beta\lambda_1\lambda_2\lambda_3\lambda_4,
\]

\MathBookSourcePage{163}
其中
\[
\alpha_i=
\left[\int_{F_i}(\mathbf v-I\mathbf v)\cdot\mathbf n_i\,\mathrm ds\right]
\left/\left[\int_{F_i}\lambda_j\lambda_k\lambda_l\,\mathrm ds\right]\right.,
\]
且
\[
\boldsymbol\beta=-4\sum_{i=1}^4\alpha_i\mathbf n_i.
\]
与 Lemma~\ref{lem:primitive-bernardi-raugel-fortin-estimates} 中一样, 有
\[
|\mathbf p_i|_{m,\kappa}
\leq C_2|\det(B_\kappa)|^{1/2}\|B_\kappa^{-1}\|^m,
\]
对 $|\lambda_1\lambda_2\lambda_3\lambda_4|_{m,\kappa}$ 也有类似表达式.
同样可推出
\[
|\alpha_i|
\leq C_3|\det(B_\kappa)|^{-1/2}\|B_\kappa\|^k
|\mathbf v|_{k,\kappa},
\qquad k=2,3.
\]
因此
\[
\left|\sum_{i=1}^4\alpha_i\mathbf p_i
+\boldsymbol\beta\lambda_1\lambda_2\lambda_3\lambda_4\right|_{m,\kappa}
\leq C_4\|B_\kappa\|^k\|B_\kappa^{-1}\|^m
|\mathbf v|_{k,\kappa}.
\]
于是 \eqref{eq:primitive-tetrahedron-second-affine-interpolation} 和
$\mathcal T_h$ 的正则性给出
\[
|r_\kappa\mathbf v-\mathbf v|_{m,\kappa}
\leq C_5\sigma^m h^{k-m}|\mathbf v|_{k,\kappa}.
\]
\end{Proof}

如前所述, 无需直接建立 inf-sup 条件, 因为 $\bar X_h\subset X_h$,
且空间对 $(\bar X_h,\bar M_h)$ 满足一致 inf-sup 条件. 只需表明
$(X_h,M_h)$ 满足适当的局部条件. 事实上, 容易验证
Theorem~\ref{thm:primitive-triangle-local-inf-sup} 的陈述和证明,
对如下局部空间对无需任何改动便成立:
\[
X_h(\kappa)=\{\,\mathbf v\in\mathcal P_2(\kappa);
\ \mathbf v|_{\partial\kappa}=\mathbf0\,\},
\qquad
M_h(\kappa)=P_1\cap L_0^2(\kappa).
\]
(注意, 证明的关键是对所有 $q\in P_1$, 函数
$\lambda_1\lambda_2\lambda_3\lambda_4\operatorname{grad}q$
属于 $X_h(\kappa)$.) 根据
Theorem~\ref{thm:primitive-local-to-global-inf-sup}, 这蕴含整体 inf-sup 条件
\eqref{eq:primitive-discrete-inf-sup}.

\setcounter{statement}{10}
\begin{Lemma}\label{lem:primitive-tetrahedron-second-global-inf-sup}
设 $\mathcal T_h$ 是 $\bar\Omega$ 的正则三角剖分,
空间 $X_h,M_h$ 由
\eqref{eq:primitive-tetrahedron-second-velocity-spaces} 和
\eqref{eq:primitive-tetrahedron-second-pressure-spaces} 定义.
则存在与 $h$ 无关的常数 $\beta^*>0$, 使得
\[
\sup_{\mathbf v_h\in X_h}
\left\{
\left|\int_\Omega q_h\operatorname{div}\mathbf v_h\,\mathrm dx\right|
/|\mathbf v_h|_{1,\Omega}
\right\}
\geq\beta^*\|q_h\|_{0,\Omega}
\qquad\forall q_h\in M_h.
\]
\end{Lemma}

最后, Hypothesis H2 约化为 $Q_h$ 上 $L^2$ 投影算子 $\rho_h$
的标准逼近性质(参见 Lemma~\ref{lem:appendix-local-l2-projection-error}):
\[
\|q-\rho_hq\|_{0,\Omega}
\leq Ch^k|q|_{k,\Omega}
\qquad\forall q\in H^k(\Omega),
\quad k=1,2.
\]
因此已经得到二阶格式所需的估计.

\MathBookSourcePage{164}
\setcounter{statement}{4}
\begin{Theorem}\label{thm:primitive-tetrahedron-second-convergence}
设 $\Omega$ 和 $\mathcal T_h$ 如
Theorem~\ref{thm:primitive-tetrahedron-first-convergence} 中所述,
Stokes 问题的解 $(\mathbf u,p)$ 满足
\[
\mathbf u\in[H^{k+1}(\Omega)\cap H_0^1(\Omega)]^3,
\qquad p\in H^k(\Omega)\cap L_0^2(\Omega),
\qquad k=1\text{ 或 }2.
\]
则采用分别由 \eqref{eq:primitive-tetrahedron-second-velocity-spaces} 和
\eqref{eq:primitive-tetrahedron-second-pressure-spaces} 定义的空间
$X_h,M_h$ 时, 格式 \eqref{eq:primitive-stokes-discrete-mixed} 的解
$(\mathbf u_h,p_h)$ 满足
\begin{equation}\label{eq:primitive-tetrahedron-second-h1-error}
|\mathbf u-\mathbf u_h|_{1,\Omega}
+\|p-p_h\|_{0,\Omega}
\leq C_1h^k(|\mathbf u|_{k+1,\Omega}+|p|_{k,\Omega}),
\qquad k=1,2.
\end{equation}
此外, 若 Stokes 问题
\eqref{eq:primitive-stokes-dual-regularity-problem} 正则, 则有
$L^2$ 估计
\begin{equation}\label{eq:primitive-tetrahedron-second-l2-error}
\|\mathbf u-\mathbf u_h\|_{0,\Omega}
\leq C_2h^{k+1}(|\mathbf u|_{k+1,\Omega}+|p|_{k,\Omega}),
\qquad k=1,2.
\end{equation}
\end{Theorem}

$l\geq3$ 次高阶格式是二阶格式的简单推广. 若要使速度空间与
$l-1$ 次多项式压力 $q$ 相配, 可以在 $P_{l+2}^3$ 的如下子空间中
选取速度 $\mathbf w$:
\begin{equation}\label{eq:primitive-tetrahedron-higher-local-space}
\mathcal P_l(\kappa)
=\left[P_l\oplus
\{\lambda_1\lambda_2\lambda_3\lambda_4
(\widetilde P_{l-2}\oplus\widetilde P_{l-3})\}\right]^3.
\end{equation}
于是 $\mathbf w$ 的每个分量在 $\kappa$ 的面上约化为 $l$ 次多项式;
此外, 只要 $q\in P_{l-1}$, 就有
$\lambda_1\lambda_2\lambda_3\lambda_4\operatorname{grad}q
\in\mathcal P_l(\kappa)$(参见
Remark~\ref{rem:primitive-bubble-gradient-property}). 因而选择如下空间:
\begin{equation}\label{eq:primitive-tetrahedron-higher-velocity-spaces}
\left\{
\begin{aligned}
W_h&=\{\,\mathbf w\in\mathcal C^0(\bar\Omega)^3;
\ \mathbf w|_\kappa\in\mathcal P_l(\kappa)
\quad\forall\kappa\in\mathcal T_h\,\},\\
X_h&=W_h\cap H_0^1(\Omega)^3,
\end{aligned}
\right.
\end{equation}
\begin{equation}\label{eq:primitive-tetrahedron-higher-pressure-spaces}
\left\{
\begin{aligned}
Q_h&=\{\,q\in L^2(\Omega);
\ q|_\kappa\in P_{l-1}\quad\forall\kappa\in\mathcal T_h\,\},\\
M_h&=Q_h\cap L_0^2(\Omega).
\end{aligned}
\right.
\end{equation}
与二维情形一样, 可以取速度 $\mathbf w$ 在 $\kappa$ 每个面 $F_i$ 上
与 $P_l$ 对应的值, 以及它在 $\kappa$ 中心处与 $P_{l-2}$ 对应的导数:
\begin{subequations}\label{eq:primitive-tetrahedron-higher-dofs}
\begin{align}
&\mathbf w(a),
&&a\in\Sigma_\kappa\cap F_i,\quad1\leq i\leq4,
\label{eq:primitive-tetrahedron-higher-dofs-boundary}\\
&\partial^k w(a_\kappa),
&&\text{$w$ 的所有 $k$ 阶偏导数},\quad0\leq k\leq l-2,
\label{eq:primitive-tetrahedron-higher-dofs-derivatives}
\end{align}
\end{subequations}
其中 $w$ 表示 $\mathbf w$ 的任意一个分量. 必要时也可以用内部矩替代导数:
\begin{equation*}\tag{2.49c}
\label{eq:primitive-tetrahedron-higher-dofs-moments}
\int_\kappa wq\,\mathrm dx
\qquad\forall q\in P_{l-2}.
\end{equation*}
一方面, 每种 \eqref{eq:primitive-tetrahedron-higher-dofs} 共定义
\[
\begin{aligned}
&4+6\dim(P_{l-2}(\mathbb R))
+4\dim(P_{l-3}(\mathbb R^2))
+\dim(P_{l-2}(\mathbb R^3))\\
&\qquad
=4+6(l-1)+2(l-2)(l-1)+\frac16(l-1)l(l+1)\\
&\qquad
=\frac16\{(l^2-1)(l+12)+24\}
\end{aligned}
\]
个自由度.

\MathBookSourcePage{165}
另一方面, $\mathcal P_l(\kappa)$ 中每个分量 $p$ 所属空间的维数为
\[
\begin{aligned}
&\dim(P_l)+\dim(\widetilde P_{l-2})+\dim(\widetilde P_{l-3}),
\qquad\text{均在 }\mathbb R^3\text{ 中},\\
&\quad=\frac16(l+1)(l+2)(l+3)
+\frac12(l-1)l+\frac12(l-2)(l-1)\\
&\quad=\frac16\{(l^2-1)(l+12)+24\}.
\end{aligned}
\]
回顾 Lemma~\ref{lem:primitive-triangle-higher-unisolvence} 的论证,
不难证明这些自由度对 $\mathcal P_l(\kappa)$ 唯一可解.

\setcounter{statement}{11}
\begin{Lemma}\label{lem:primitive-tetrahedron-higher-unisolvence}
$\mathcal P_l(\kappa)$ 中多项式的每个分量 $p$ 由自由度
\eqref{eq:primitive-tetrahedron-higher-dofs-boundary} 和
\eqref{eq:primitive-tetrahedron-higher-dofs-derivatives}, 或者
\eqref{eq:primitive-tetrahedron-higher-dofs-boundary} 和
\eqref{eq:primitive-tetrahedron-higher-dofs-moments} 唯一确定.
此外, 在 $\kappa$ 的给定面上, $p$ 只依赖于定义在该面上的自由度.
\end{Lemma}

该 Lemma 给出如下插值算子:
\[
r_h\mathbf v|_\kappa=r_\kappa\mathbf v
\qquad\forall\kappa\in\mathcal T_h,
\quad\forall\mathbf v\in\mathcal C^0(\bar\Omega)^3,
\]
其中 $r_\kappa\mathbf v$ 是 $\mathcal P_l(\kappa)$ 中与 $\mathbf v$
具有相同自由度
\eqref{eq:primitive-tetrahedron-higher-dofs-boundary} 和
\eqref{eq:primitive-tetrahedron-higher-dofs-moments} 的唯一多项式.
显然, 算子 $r_\kappa$ 在仿射变换下不变, 因而有如下结果.

\setcounter{statement}{12}
\begin{Lemma}\label{lem:primitive-tetrahedron-higher-interpolation}
算子 $r_h$ 属于
\[
\mathcal L(\mathcal C^0(\bar\Omega)^3;W_h)
\cap\mathcal L([\mathcal C^0(\bar\Omega)\cap H_0^1(\Omega)]^3;X_h).
\]
此外, 若三角剖分 $\mathcal T_h$ 正则, 则 $r_h$ 具有如下插值误差:
\begin{equation}\label{eq:primitive-tetrahedron-higher-interpolation}
|\mathbf v-r_h\mathbf v|_{m,\Omega}
\leq Ch^{k+1-m}|\mathbf v|_{k+1,\Omega}
\qquad\forall\mathbf v\in H^{k+1}(\Omega)^3,
\quad m=0\text{ 或 }1,
\end{equation}
其中 $1\leq k\leq l$, 正常数 $C$ 与 $h$ 和 $\mathbf v$ 无关.
\end{Lemma}

与二阶格式一样, 此处 inf-sup 条件的证明与
Theorem~\ref{thm:primitive-triangle-local-inf-sup} 的证明完全类似.
所涉及的局部空间为
\[
X_h(\kappa)=\{\,\mathbf v\in\mathcal P_l(\kappa);
\ \mathbf v|_{\partial\kappa}=\mathbf0\,\},
\qquad
M_h(\kappa)=P_{l-1}\cap L_0^2(\Omega),
\]
将它们联系起来的关键性质是
\[
\lambda_1\lambda_2\lambda_3\lambda_4\operatorname{grad}q
\in X_h(\kappa)
\qquad\forall q\in M_h(\kappa).
\]
因此, Lemma~\ref{lem:primitive-tetrahedron-second-global-inf-sup}
的陈述同样适用于由
\eqref{eq:primitive-tetrahedron-higher-velocity-spaces} 和
\eqref{eq:primitive-tetrahedron-higher-pressure-spaces} 定义的空间对
$(X_h,M_h)$.

Hypothesis H2 仍由 $Q_h$ 上 $L^2$ 投影算子 $\rho_h$ 的逼近性质得到:
\[
\|q-\rho_hq\|_{0,\Omega}
\leq Ch^k|q|_{k,\Omega}
\qquad\forall q\in H^k(\Omega),
\quad1\leq k\leq l.
\]
汇集这些结果, 即得到该最后一个格式所需的估计.
